\IfFileExists{stacks-project.cls}{%
\documentclass{stacks-project}
}{%
\documentclass{amsart}
}

% For dealing with references we use the comment environment
\usepackage{verbatim}
\newenvironment{reference}{\comment}{\endcomment}
%\newenvironment{reference}{}{}
\newenvironment{slogan}{\comment}{\endcomment}
\newenvironment{history}{\comment}{\endcomment}

% For commutative diagrams we use Xy-pic
\usepackage[all]{xy}

% We use 2cell for 2-commutative diagrams.
\xyoption{2cell}
\UseAllTwocells

% We use multicol for the list of chapters between chapters
\usepackage{multicol}

% This is generally recommended for better output
\usepackage{lmodern}
\usepackage[T1]{fontenc}

% For cross-file-references
\usepackage{xr-hyper}

% Package for hypertext links:
\usepackage{hyperref}

% For any local file, say "hello.tex" you want to link to please
% use \externaldocument[hello-]{hello}
\externaldocument[introduction-]{introduction}
\externaldocument[conventions-]{conventions}
\externaldocument[sets-]{sets}
\externaldocument[categories-]{categories}
\externaldocument[topology-]{topology}
\externaldocument[sheaves-]{sheaves}
\externaldocument[sites-]{sites}
\externaldocument[stacks-]{stacks}
\externaldocument[fields-]{fields}
\externaldocument[algebra-]{algebra}
\externaldocument[brauer-]{brauer}
\externaldocument[homology-]{homology}
\externaldocument[derived-]{derived}
\externaldocument[simplicial-]{simplicial}
\externaldocument[more-algebra-]{more-algebra}
\externaldocument[smoothing-]{smoothing}
\externaldocument[modules-]{modules}
\externaldocument[sites-modules-]{sites-modules}
\externaldocument[injectives-]{injectives}
\externaldocument[cohomology-]{cohomology}
\externaldocument[sites-cohomology-]{sites-cohomology}
\externaldocument[dga-]{dga}
\externaldocument[dpa-]{dpa}
\externaldocument[sdga-]{sdga}
\externaldocument[hypercovering-]{hypercovering}
\externaldocument[schemes-]{schemes}
\externaldocument[constructions-]{constructions}
\externaldocument[properties-]{properties}
\externaldocument[morphisms-]{morphisms}
\externaldocument[coherent-]{coherent}
\externaldocument[divisors-]{divisors}
\externaldocument[limits-]{limits}
\externaldocument[varieties-]{varieties}
\externaldocument[topologies-]{topologies}
\externaldocument[descent-]{descent}
\externaldocument[perfect-]{perfect}
\externaldocument[more-morphisms-]{more-morphisms}
\externaldocument[flat-]{flat}
\externaldocument[groupoids-]{groupoids}
\externaldocument[more-groupoids-]{more-groupoids}
\externaldocument[etale-]{etale}
\externaldocument[chow-]{chow}
\externaldocument[intersection-]{intersection}
\externaldocument[pic-]{pic}
\externaldocument[weil-]{weil}
\externaldocument[adequate-]{adequate}
\externaldocument[dualizing-]{dualizing}
\externaldocument[duality-]{duality}
\externaldocument[discriminant-]{discriminant}
\externaldocument[derham-]{derham}
\externaldocument[local-cohomology-]{local-cohomology}
\externaldocument[algebraization-]{algebraization}
\externaldocument[curves-]{curves}
\externaldocument[resolve-]{resolve}
\externaldocument[models-]{models}
\externaldocument[functors-]{functors}
\externaldocument[equiv-]{equiv}
\externaldocument[pione-]{pione}
\externaldocument[etale-cohomology-]{etale-cohomology}
\externaldocument[proetale-]{proetale}
\externaldocument[relative-cycles-]{relative-cycles}
\externaldocument[more-etale-]{more-etale}
\externaldocument[trace-]{trace}
\externaldocument[crystalline-]{crystalline}
\externaldocument[spaces-]{spaces}
\externaldocument[spaces-properties-]{spaces-properties}
\externaldocument[spaces-morphisms-]{spaces-morphisms}
\externaldocument[decent-spaces-]{decent-spaces}
\externaldocument[spaces-cohomology-]{spaces-cohomology}
\externaldocument[spaces-limits-]{spaces-limits}
\externaldocument[spaces-divisors-]{spaces-divisors}
\externaldocument[spaces-over-fields-]{spaces-over-fields}
\externaldocument[spaces-topologies-]{spaces-topologies}
\externaldocument[spaces-descent-]{spaces-descent}
\externaldocument[spaces-perfect-]{spaces-perfect}
\externaldocument[spaces-more-morphisms-]{spaces-more-morphisms}
\externaldocument[spaces-flat-]{spaces-flat}
\externaldocument[spaces-groupoids-]{spaces-groupoids}
\externaldocument[spaces-more-groupoids-]{spaces-more-groupoids}
\externaldocument[bootstrap-]{bootstrap}
\externaldocument[spaces-pushouts-]{spaces-pushouts}
\externaldocument[spaces-chow-]{spaces-chow}
\externaldocument[groupoids-quotients-]{groupoids-quotients}
\externaldocument[spaces-more-cohomology-]{spaces-more-cohomology}
\externaldocument[spaces-simplicial-]{spaces-simplicial}
\externaldocument[spaces-duality-]{spaces-duality}
\externaldocument[formal-spaces-]{formal-spaces}
\externaldocument[restricted-]{restricted}
\externaldocument[spaces-resolve-]{spaces-resolve}
\externaldocument[formal-defos-]{formal-defos}
\externaldocument[defos-]{defos}
\externaldocument[cotangent-]{cotangent}
\externaldocument[examples-defos-]{examples-defos}
\externaldocument[algebraic-]{algebraic}
\externaldocument[examples-stacks-]{examples-stacks}
\externaldocument[stacks-sheaves-]{stacks-sheaves}
\externaldocument[criteria-]{criteria}
\externaldocument[artin-]{artin}
\externaldocument[quot-]{quot}
\externaldocument[stacks-properties-]{stacks-properties}
\externaldocument[stacks-morphisms-]{stacks-morphisms}
\externaldocument[stacks-limits-]{stacks-limits}
\externaldocument[stacks-cohomology-]{stacks-cohomology}
\externaldocument[stacks-perfect-]{stacks-perfect}
\externaldocument[stacks-introduction-]{stacks-introduction}
\externaldocument[stacks-more-morphisms-]{stacks-more-morphisms}
\externaldocument[stacks-geometry-]{stacks-geometry}
\externaldocument[moduli-]{moduli}
\externaldocument[moduli-curves-]{moduli-curves}
\externaldocument[examples-]{examples}
\externaldocument[exercises-]{exercises}
\externaldocument[guide-]{guide}
\externaldocument[desirables-]{desirables}
\externaldocument[coding-]{coding}
\externaldocument[obsolete-]{obsolete}
\externaldocument[fdl-]{fdl}
\externaldocument[index-]{index}

% Theorem environments.
%
\theoremstyle{plain}
\newtheorem{theorem}[subsection]{Theorem}
\newtheorem{proposition}[subsection]{Proposition}
\newtheorem{lemma}[subsection]{Lemma}

\theoremstyle{definition}
\newtheorem{definition}[subsection]{Definition}
\newtheorem{example}[subsection]{Example}
\newtheorem{exercise}[subsection]{Exercise}
\newtheorem{situation}[subsection]{Situation}

\theoremstyle{remark}
\newtheorem{remark}[subsection]{Remark}
\newtheorem{remarks}[subsection]{Remarks}

\numberwithin{equation}{subsection}

% Macros
%
\def\lim{\mathop{\mathrm{lim}}\nolimits}
\def\colim{\mathop{\mathrm{colim}}\nolimits}
\def\Spec{\mathop{\mathrm{Spec}}}
\def\Hom{\mathop{\mathrm{Hom}}\nolimits}
\def\Ext{\mathop{\mathrm{Ext}}\nolimits}
\def\SheafHom{\mathop{\mathcal{H}\!\mathit{om}}\nolimits}
\def\SheafExt{\mathop{\mathcal{E}\!\mathit{xt}}\nolimits}
\def\Sch{\mathit{Sch}}
\def\Mor{\mathop{\mathrm{Mor}}\nolimits}
\def\Ob{\mathop{\mathrm{Ob}}\nolimits}
\def\Sh{\mathop{\mathit{Sh}}\nolimits}
\def\NL{\mathop{N\!L}\nolimits}
\def\CH{\mathop{\mathrm{CH}}\nolimits}
\def\proetale{{pro\text{-}\acute{e}tale}}
\def\etale{{\acute{e}tale}}
\def\QCoh{\mathit{QCoh}}
\def\Ker{\mathop{\mathrm{Ker}}}
\def\Im{\mathop{\mathrm{Im}}}
\def\Coker{\mathop{\mathrm{Coker}}}
\def\Coim{\mathop{\mathrm{Coim}}}

% Boxtimes
%
\DeclareMathSymbol{\boxtimes}{\mathbin}{AMSa}{"02}

%
% Macros for moduli stacks/spaces
%
\def\QCohstack{\mathcal{QC}\!\mathit{oh}}
\def\Cohstack{\mathcal{C}\!\mathit{oh}}
\def\Spacesstack{\mathcal{S}\!\mathit{paces}}
\def\Quotfunctor{\mathrm{Quot}}
\def\Hilbfunctor{\mathrm{Hilb}}
\def\Curvesstack{\mathcal{C}\!\mathit{urves}}
\def\Polarizedstack{\mathcal{P}\!\mathit{olarized}}
\def\Complexesstack{\mathcal{C}\!\mathit{omplexes}}
% \Pic is the operator that assigns to X its picard group, usage \Pic(X)
% \Picardstack_{X/B} denotes the Picard stack of X over B
% \Picardfunctor_{X/B} denotes the Picard functor of X over B
\def\Pic{\mathop{\mathrm{Pic}}\nolimits}
\def\Picardstack{\mathcal{P}\!\mathit{ic}}
\def\Picardfunctor{\mathrm{Pic}}
\def\Deformationcategory{\mathcal{D}\!\mathit{ef}}
\usepackage{cleveref}
\crefname{originallemma}{original lemma}{original lemmas}
\Crefname{originallemma}{Original lemma}{Original lemmas}
\makeatletter
\newlabel{lemma-lqf-shriek-derived}{{0F5N}{source}{}{sourcepointer-1}{}}
\newlabel{lemma-lqf-shriek-derived@cref}{{[originallemma][0][]0F5N}{[1][0][]source}{}{}{}}
\makeatother
\crefname{originallemma}{original lemma}{original lemmas}
\Crefname{originallemma}{Original lemma}{Original lemmas}
\makeatletter
\newlabel{lemma-lqf-colimit-f-shriek}{{0F5H}{source}{}{sourcepointer-2}{}}
\newlabel{lemma-lqf-colimit-f-shriek@cref}{{[originallemma][0][]0F5H}{[1][0][]source}{}{}{}}
\makeatother
\crefname{originallemma}{original lemma}{original lemmas}
\Crefname{originallemma}{Original lemma}{Original lemmas}
\makeatletter
\newlabel{lemma-f-shriek-composition}{{0F50}{source}{}{sourcepointer-3}{}}
\newlabel{lemma-f-shriek-composition@cref}{{[originallemma][0][]0F50}{[1][0][]source}{}{}{}}
\makeatother
\crefname{originallemma}{original lemma}{original lemmas}
\Crefname{originallemma}{Original lemma}{Original lemmas}
\makeatletter
\newlabel{lemma-proper-f-shriek}{{0F51}{source}{}{sourcepointer-4}{}}
\newlabel{lemma-proper-f-shriek@cref}{{[originallemma][0][]0F51}{[1][0][]source}{}{}{}}
\makeatother
\crefname{originallemma}{original lemma}{original lemmas}
\Crefname{originallemma}{Original lemma}{Original lemmas}
\makeatletter
\newlabel{more-morphisms-lemma-quasi-finite-separated-pass-through-finite-addendum}{{0F2N}{source}{}{sourcepointer-5}{}}
\newlabel{more-morphisms-lemma-quasi-finite-separated-pass-through-finite-addendum@cref}{{[originallemma][0][]0F2N}{[1][0][]source}{}{}{}}
\makeatother
\crefname{originaldefinition}{original definition}{original definitions}
\Crefname{originaldefinition}{Original definition}{Original definitions}
\makeatletter
\newlabel{morphisms-definition-finite-presentation}{{01TP}{source}{}{sourcepointer-6}{}}
\newlabel{morphisms-definition-finite-presentation@cref}{{[originaldefinition][0][]01TP}{[1][0][]source}{}{}{}}
\makeatother
\crefname{originallemma}{original lemma}{original lemmas}
\Crefname{originallemma}{Original lemma}{Original lemmas}
\makeatletter
\newlabel{etale-cohomology-lemma-jshriek-constructible}{{03S8}{source}{}{sourcepointer-7}{}}
\newlabel{etale-cohomology-lemma-jshriek-constructible@cref}{{[originallemma][0][]03S8}{[1][0][]source}{}{}{}}
\makeatother
\crefname{originallemma}{original lemma}{original lemmas}
\Crefname{originallemma}{Original lemma}{Original lemmas}
\makeatletter
\newlabel{etale-cohomology-lemma-finite-pushforward-constructible}{{095R}{source}{}{sourcepointer-8}{}}
\newlabel{etale-cohomology-lemma-finite-pushforward-constructible@cref}{{[originallemma][0][]095R}{[1][0][]source}{}{}{}}
\makeatother
\title{Lower shriek for locally quasi-finite finitely presented morphisms preserves constructibility}
\author{The Stacks Project Authors}
\date{stacks-project:more-etale.tex:local-snapshot-20261012; GNU FDL 1.2 or later}
\begin{document}
\maketitle
\setcounter{subsection}{1}
\noindent\textit{Editorial note.} All delivered mathematical text is the author's own, copied from the pinned Stacks Project chapter file more-etale.tex; no statement or proof step was rewritten, reordered or removed. The counted span is the author environment and its own proof environment. References inside the span are handed to the assembler as navigation stubs that carry the canonical tag of the cited tag only; whether each of them is load-bearing is recorded on the stub itself.\par\medskip
\noindent\textit{Allowed standard prerequisites (editorial).}\begin{itemize}
\item Construction of the lower shriek functor on derived categories for locally quasi-finite morphisms (Lemma lemma-lqf-shriek-derived, canonical tag 0F5N). cited at statement lines 4690-4691. The functor f\_! whose constructibility preservation is claimed is the functor constructed there.
\item Colimit description of lower shriek (Lemma lemma-lqf-colimit-f-shriek, canonical tag 0F5H). cited at proof lines 4704-4708. Reduces the claim to the pieces f\_\{i,!\} and f\_\{ii',!\} over the finite affine covering.
\item Zariski main theorem factorization through a finite morphism (More on Morphisms, Lemma more-morphisms-lemma-quasi-finite-separated-pass-through-finite-addendum, canonical tag 0F2N). cited at proof lines 4711-4716. Supplies the factorization f = g after j with j an open immersion and g finite of finite presentation.
\item Composition of lower shriek (Lemma lemma-f-shriek-composition, canonical tag 0F50). cited at proof lines 4716-4717. Gives f\_! = g\_! after j\_! for the factorization.
\item Constructibility of lower shriek along open immersions (Etale Cohomology, Lemma etale-cohomology-lemma-jshriek-constructible, canonical tag 03S8). cited at proof lines 4718-4719. j\_! of the constructible sheaf F is constructible on X'.
\item Agreement of lower shriek with pushforward for proper morphisms (Lemma lemma-proper-f-shriek, canonical tag 0F51). cited at proof lines 4720-4721. The finite morphism g satisfies g\_! = g\_*.
\item Constructibility of pushforward along finite morphisms (Etale Cohomology, Lemma etale-cohomology-lemma-finite-pushforward-constructible, canonical tag 095R). cited at proof lines 4722-4725. g\_* of the constructible sheaf is constructible, which closes the argument.
\end{itemize}\medskip
\ifdefined\hypertarget\hypertarget{sourcepointer-1}{}\fi
\noindent\textit{Source reference.} Original lemma 0F5N: sources/stacks-chow/tags.txt, source lines 18943--18943. Cited at statement lines 4690-4691: the functor f\_! whose constructibility preservation is claimed is the functor constructed there.\par\medskip
\ifdefined\hypertarget\hypertarget{sourcepointer-2}{}\fi
\noindent\textit{Source reference.} Original lemma 0F5H: sources/stacks-chow/tags.txt, source lines 18937--18937. Cited at proof lines 4704-4708: reduces the claim to the pieces f\_\{i,!\} and f\_\{ii',!\} attached to the finite affine covering.\par\medskip
\ifdefined\hypertarget\hypertarget{sourcepointer-3}{}\fi
\noindent\textit{Source reference.} Original lemma 0F50: sources/stacks-chow/tags.txt, source lines 18920--18920. Cited at proof lines 4716-4717: gives f\_! = g\_! \textbackslash{}circ j\_! for the factorization f = g \textbackslash{}circ j.\par\medskip
\ifdefined\hypertarget\hypertarget{sourcepointer-4}{}\fi
\noindent\textit{Source reference.} Original lemma 0F51: sources/stacks-chow/tags.txt, source lines 18921--18921. Cited at proof lines 4720-4721: the finite morphism g satisfies g\_! = g\_*.\par\medskip
\ifdefined\hypertarget\hypertarget{sourcepointer-5}{}\fi
\noindent\textit{Source reference.} Original lemma 0F2N: sources/stacks-chow/tags.txt, source lines 18838--18838. Cited at proof lines 4711-4716: supplies the factorization f = g \textbackslash{}circ j with j an open immersion and g finite of finite presentation.\par\medskip
\ifdefined\hypertarget\hypertarget{sourcepointer-6}{}\fi
\noindent\textit{Source reference.} Original definition 01TP: sources/stacks-chow/tags.txt, source lines 2382--2382. Cited at proof lines 4701-4702: X is quasi-compact and quasi-separated when f is locally quasi-finite and of finite presentation.\par\medskip
\ifdefined\hypertarget\hypertarget{sourcepointer-7}{}\fi
\noindent\textit{Source reference.} Original lemma 03S8: sources/stacks-chow/tags.txt, source lines 4957--4957. Cited at proof lines 4718-4719: j\_!F is constructible on X'.\par\medskip
\ifdefined\hypertarget\hypertarget{sourcepointer-8}{}\fi
\noindent\textit{Source reference.} Original lemma 095R: sources/stacks-chow/tags.txt, source lines 11594--11594. Cited at proof lines 4722-4725: g\_*j\_!F is constructible, which closes the argument.\par\medskip
\begin{lemma}
\label{lemma-qf-f-shriek-constructible}
Let $f : X \to Y$ be a morphism of schemes which is
locally quasi-finite and of finite presentation.
The functor $f_! : D(X_\etale, \Lambda) \to D(Y_\etale, \Lambda)$
of Lemma \ref{lemma-lqf-shriek-derived}
sends $D_c(X_\etale, \Lambda)$ into $D_c(Y_\etale, \Lambda)$.
\end{lemma}

\begin{proof}
Since the functor $f_!$ is exact, it suffices to show that
$f_!\mathcal{F}$ is constructible for any constructible sheaf
$\mathcal{F}$ of $\Lambda$-modules on $X_\etale$.
The question is local on $Y$ and hence
we may and do assume $Y$ is affine.
Then $X$ is quasi-compact and quasi-separated, see
Morphisms, Definition \ref{morphisms-definition-finite-presentation}.
Say $X = \bigcup_{i = 1, \ldots, n} X_i$ is a finite affine open
covering. By Lemma \ref{lemma-lqf-colimit-f-shriek}
we see that it suffices to show that
$f_{i, !}\mathcal{F}|_{X_i}$ and $f_{ii', !}\mathcal{F}|_{X_i \cap X_{i'}}$
are constructible where $f_i : X_i \to Y$ and
$f_{ii'} : X_i \cap X_{i'} \to Y$ are the restrictions of $f$.
Since $X_i$ and $X_i \cap X_{i'}$ are quasi-compact and separated
this means we may assume $f$ is separated.
By Zariski's main theorem (in the form of
More on Morphisms, Lemma
\ref{more-morphisms-lemma-quasi-finite-separated-pass-through-finite-addendum})
we can choose a factorization $f = g \circ j$ where
$j : X \to X'$ is an open immersion and $g : X' \to Y$ is
finite and of finite presentation. Then $f_! = g_! \circ j_!$ by
Lemma \ref{lemma-f-shriek-composition}.
By \'Etale Cohomology, Lemma \ref{etale-cohomology-lemma-jshriek-constructible}
we see that $j_!\mathcal{F}$ is constructible on $X'$.
The morphism $g$ is finite hence $g_! = g_*$
by Lemma \ref{lemma-proper-f-shriek}.
Thus $f_!\mathcal{F} = g_!j_!\mathcal{F} = g_*j_!\mathcal{F}$
is constructible by
\'Etale Cohomology, Lemma
\ref{etale-cohomology-lemma-finite-pushforward-constructible}.
\end{proof}

\end{document}
